\documentclass[12pt]{article}

\usepackage[utf8]{inputenc}
\usepackage[T1]{fontenc}
\usepackage{setspace}
\usepackage[margin=1in]{geometry}
\usepackage{amsmath,amssymb}
\usepackage{hyperref}
\usepackage{lmodern}
\usepackage{microtype}

\hypersetup{colorlinks=true, linkcolor=blue, citecolor=blue, urlcolor=blue}

\title{Fraud Allegations as Identification:\\
       A Within-Election Forensic Test Using Institutional Challenge Records}

\author{Carlos C\'esar Ch\'avez Padilla\thanks{%
  Email: \href{mailto:cchavezp@uchicago.edu}{cchavezp@uchicago.edu}.
  I thank the Oficina Nacional de Procesos Electorales for making
  mesa-level data publicly available. All errors are my own.
  \textit{Data availability:} The primary data are publicly available
  at \url{https://datosabiertos.gob.pe}. Replication code is available
  at \url{https://github.com/cesarchavezp29/peru-election-forensics-2021}.}\\[4pt]
  \normalsize Center for the Economics of Human Development\\
  \normalsize University of Chicago}

\date{September 2024}

\begin{document}

\maketitle
\thispagestyle{empty}

\begin{abstract}
\noindent
\begin{onehalfspacing}
\small
Standard election forensics methods---distributional fingerprint analysis,
digit tests, ecological regression---apply cross-sectional diagnostics that
cannot distinguish genuine geographic polarization from ballot manipulation
when both produce similar distributional signatures. This paper introduces a
within-election identification strategy that exploits institutional challenge
records: when losing candidates file formal nullification challenges against
specific polling stations, they inadvertently define a treatment group whose
forensic properties can be compared to the unchallenged remainder of the same
election. The design uses the fraud allegations themselves to construct the
counterfactual, turning the challenger's own legal record into an
identification instrument. I apply this strategy to Peru's 2021 presidential
runoff---the narrowest result in modern Peruvian history, followed by six
weeks of fraud allegations and delayed certification---using the complete
universe of 82{,}870 polling-station records. No forensic method detects
manipulation. Contested polling stations, far from over-performing for the
alleged beneficiary, record a statistically significant two-percentage-point
lower swing toward Castillo than comparable uncontested stations. The paper
also documents how extreme geographic polarization generates distributional
patterns observationally similar to ballot-stuffing signatures, with
implications for the application of fingerprint methods in highly sorted
electorates.
\end{onehalfspacing}

\medskip
\noindent\textbf{Keywords:} election forensics, electoral fraud, Peru,
geographic polarization, distributional fingerprints, digit tests.

\medskip
\noindent\textbf{JEL codes:} D72, C12, P16.
\end{abstract}

\thispagestyle{empty}
\end{document}
